% chapters/17/section-03.tex
\section{Satu Permainan, Banyak Kompetensi}
\label{sec:lapisan-kompetensi}

\subsection{Mengapa Menguasai Setiap Bab Belum Cukup}

Seluruh bab pada buku ini dapat Anda tuntaskan satu per satu, dan tetap
menghasilkan permainan yang tidak dapat dimainkan. Itu bukan kemungkinan yang
jauh: ia adalah hasil yang paling sering terjadi pada proyek kelompok yang
dikerjakan secara terpisah-pisah.

Sebabnya ada pada cara kompetensi itu disusun. Setiap bab melatih satu
kemampuan pada keadaan yang sudah disiapkan. Bab~VII melatih tabrakan pada dua
benda; Bab~XIV melatih animasi pada satu pemain; Bab~XV melatih bunyi pada satu
peristiwa. Ketika seluruhnya digabungkan, keadaan yang Anda hadapi berbeda:
banyak benda bergerak sekaligus, peristiwa terjadi bersamaan, dan kesalahan
pada satu bagian tampak sebagai kesalahan pada bagian lain.

Gambar~\ref{fig:lapisan-permainan} memperlihatkan lapisan-lapisan yang menyusun
sebuah permainan yang selesai, beserta bab tempat setiap lapisan dibangun.

\begin{figure}[htbp]
\centering
% chapters/17/figures\lapisan-permainan.tikz
% Lapisan-lapisan yang menyusun satu permainan yang selesai, beserta bab
% tempat setiap lapisan dibangun.
\begin{tikzpicture}[
  lapisan/.style={
    rectangle, rounded corners=4pt, draw=ubPrimary, thick, fill=ubAlt,
    align=center, text width=96mm, inner sep=4pt, minimum height=8mm,
    font=\footnotesize
  },
  dasar/.style={
    rectangle, rounded corners=4pt, draw=ubAccent, thick, fill=white,
    align=center, text width=96mm, inner sep=4pt, minimum height=8mm,
    font=\footnotesize
  },
  sisi/.style={font=\scriptsize\bfseries, ubPrimary, align=left}
]

% Ditulis dari bawah ke atas: lapisan yang paling dahulu dibangun berada
% di bawah, sebab lapisan di atasnya bergantung padanya.
\node[dasar]  (l1) at (5.0,0.0)  {\textbf{Bab IV} --- \istilah{Game loop}, masukan, dan \istilah{delta time}};
\node[dasar]  (l2) at (5.0,0.95) {\textbf{Bab V} --- Gerak objek dan vektor};
\node[dasar]  (l3) at (5.0,1.9)  {\textbf{Bab VI} --- Fisika sederhana: gravitasi dan lompatan};
\node[dasar]  (l4) at (5.0,2.85) {\textbf{Bab VII} --- \istilah{Sprite}, kelompok, dan tabrakan};
\node[lapisan](l5) at (5.0,3.8)  {\textbf{Bab VIII--XI} --- Perilaku NPC, mesin keadaan, dan pencarian jalur};
\node[lapisan](l6) at (5.0,4.75) {\textbf{Bab XII--XIII} --- Level Unity, fisika mesin, dan kendali pemain};
\node[lapisan](l7) at (5.0,5.7)  {\textbf{Bab XIV} --- Aset visual, animasi, dan antarmuka};
\node[lapisan](l8) at (5.0,6.65) {\textbf{Bab XV} --- Bunyi dan umpan balik interaksi};
\node[lapisan](l9) at (5.0,7.6)  {\textbf{Bab XVI} --- Pengujian, perbaikan, dan penyajian};

\node[sisi, anchor=east] at (-1.4,1.42) {Dasar};
\node[sisi, anchor=east] at (-1.4,3.8)  {Kecerdasan};
\node[sisi, anchor=east] at (-1.4,5.7)  {Penyajian};
\node[sisi, anchor=east] at (-1.4,7.6)  {Penilaian};
\end{tikzpicture}

\caption{Lapisan-lapisan yang menyusun satu permainan, dan bab tempat setiap lapisan dibangun}
\label{fig:lapisan-permainan}
\end{figure}

Susunan itu bukan sekadar urutan pengerjaan. Setiap lapisan bergantung pada
lapisan di bawahnya, dan ketergantungan itulah yang menentukan arah penelusuran
ketika ada yang salah. Bunyi yang tidak terdengar pada Bab~XV dapat disebabkan
oleh peristiwanya yang tidak pernah terjadi, dan peristiwa itu tidak terjadi
karena benda yang seharusnya disentuh lolos dari pemeriksaan tabrakan pada
Bab~VII.

\subsection{Arah Penelusuran ketika Ada yang Salah}

Karena lapisan di atas bergantung pada lapisan di bawah, penelusuran selalu
bergerak ke bawah, tidak ke atas. Tiga contoh berikut memperlihatkan arah itu.

\begin{enumerate}
  \item \textbf{Skor tidak bertambah ketika benda diambil.} Jangan memeriksa
        antarmuka terlebih dahulu. Periksa apakah sentuhannya terdeteksi
        (Bab~VII), lalu apakah tanda bendanya benar, lalu apakah fungsinya
        terpanggil, dan barulah periksa tulisannya (Bab~XIV).
  \item \textbf{Pemain menembus lantai.} Jangan menyunting skrip antarmuka.
        Periksa \istilah{Rigidbody} dan \istilah{Collider} (Bab~XIII), lalu
        apakah geraknya dijalankan pada \code{FixedUpdate} (Bab~XIII), lalu
        apakah \istilah{delta time} dipakai dengan benar (Bab~IV).
  \item \textbf{Animasi tidak berganti.} Periksa dahulu apakah nilainya
        berubah (Bab~XIV), lalu apakah nilai itu berasal dari keadaan pemain
        yang benar (Bab~XIII), lalu apakah keadaan pemain itu sendiri benar
        (Bab~IV--VI).
\end{enumerate}

\begin{tipsbox}
Kebiasaan yang menghemat waktu: ketika ada yang salah, tanyakan ``lapisan
paling bawah mana yang mungkin menyebabkan gejala ini'', lalu periksa dari
bawah ke atas. Memeriksa dari atas ke bawah hampir selalu berakhir dengan
memperbaiki tampilan dari masalah yang sebenarnya ada pada fisika.
\end{tipsbox}

\subsection{Integrasi sebagai Kompetensi Tersendiri}

Kemampuan menggabungkan bagian-bagian itu sendiri adalah kompetensi, dan ia
dinilai secara terpisah. Itulah sebabnya Proyek Akhir memiliki bobot terbesar,
dan itulah sebabnya rubriknya memuat kriteria ``integrasi aset visual,
animasi, dan UI'' serta ``integrasi audio dan umpan balik interaksi'' sebagai
kriteria tersendiri, bukan sebagai bagian dari kriteria mekanik.

Pekerjaan menggabungkan itu tidak dapat dibagi menjadi tugas-tugas kecil yang
dikerjakan orang berbeda tanpa saling berbicara. Ia menuntut satu orang yang
memahami seluruh lapisan. Yang juga diperlukan adalah satu keputusan bersama
tentang hal-hal yang tampak sepele: siapa yang menyimpan skor, apa yang terjadi
ketika
nyawa habis, dan urutan apa yang harus dilalui pemain dari awal sampai tujuan.

\begin{catatabox}
Apabila kelompok Anda mengerjakan proyek secara terpisah lalu menggabungkannya
pada hari terakhir, hampir pasti ada bagian yang tidak tersambung. Cara yang
lebih baik adalah menyambungkan dua bagian setiap kali satu bagian selesai,
sebagaimana siklus pada Gambar~\ref{fig:siklus-uji} pada Bab~XVI.
\end{catatabox}
